% !TeX program = xelatex
% !TeX encoding = UTF-8
% ==============================================================================
% KHUNG LÝ THUYẾT MÔ HÌNH
% Hệ ký hiệu thống nhất với "Bài báo chính.tex"; nhãn phương trình dùng lại
% cùng tên để hai văn bản đối chiếu chéo được.
% ==============================================================================
\documentclass[12pt,a4paper]{article}

\usepackage{fontspec}
\usepackage{unicode-math}
\IfFontExistsTF{Times New Roman}{\setmainfont{Times New Roman}}{\setmainfont{Liberation Serif}}
\IfFontExistsTF{Cambria Math}{\setmathfont{Cambria Math}}{\setmathfont{LibertinusMath-Regular.otf}}

\usepackage[a4paper,top=2.5cm,bottom=2.5cm,left=2.5cm,right=2.5cm]{geometry}
\usepackage{setspace}
\onehalfspacing
\usepackage{indentfirst}
\setlength{\parindent}{0.7cm}
\setlength{\parskip}{0.2em}
\emergencystretch=2em

\usepackage{microtype}
\usepackage{amsmath}
\usepackage{amsthm}
\numberwithin{equation}{section}

\usepackage{booktabs}
\usepackage{tabularx}
\usepackage{array}
\usepackage{longtable}
\usepackage{float}
\usepackage{enumitem}
\setlist[itemize]{leftmargin=1.5em,itemsep=0.15em,topsep=0.15em}
\setlist[enumerate]{leftmargin=1.8em,itemsep=0.15em,topsep=0.15em}

\usepackage{fancyhdr}
\pagestyle{fancy}
\fancyhf{}
\fancyhead[L]{\footnotesize\itshape Khung lý thuyết MS-GARCH-X Student-t}
\fancyhead[R]{\footnotesize\thepage}
\renewcommand{\headrulewidth}{0.4pt}

\usepackage{caption}
\captionsetup{font=small,labelfont=bf}
\renewcommand{\tablename}{Bảng}
\renewcommand{\figurename}{Hình}
\renewcommand{\proofname}{Chứng minh}
\renewcommand{\abstractname}{Tóm tắt}
\renewcommand{\refname}{Tài liệu tham khảo}

\newcolumntype{Y}{>{\raggedright\arraybackslash}X}

\usepackage[hidelinks,unicode]{hyperref}
\usepackage{cleveref}
\crefname{section}{Mục}{Các mục}
\Crefname{section}{Mục}{Các mục}
\crefname{subsection}{Mục}{Các mục}
\Crefname{subsection}{Mục}{Các mục}
\crefname{equation}{}{}
\Crefname{equation}{}{}
\crefname{table}{Bảng}{Các bảng}
\Crefname{table}{Bảng}{Các bảng}
\crefname{assumption}{Giả định}{Các giả định}
\Crefname{assumption}{Giả định}{Các giả định}
\crefname{definition}{Định nghĩa}{Các định nghĩa}
\Crefname{definition}{Định nghĩa}{Các định nghĩa}

\theoremstyle{definition}
\newtheorem{definition}{Định nghĩa}[section]
\newtheorem{assumption}{Giả định}[section]
\newtheorem{proposition}{Mệnh đề}[section]
\newtheorem{remark}{Nhận xét}[section]

% --- Ký hiệu: đồng bộ với "Bài báo chính.tex" -------------------------------
\newcommand{\R}{\mathbb{R}}
\newcommand{\E}{\mathbb{E}}
\newcommand{\Var}{\operatorname{Var}}
\newcommand{\Cov}{\operatorname{Cov}}
\newcommand{\Prob}{\mathbb{P}}
\newcommand{\diag}{\operatorname{diag}}
\newcommand{\softplus}{\operatorname{softplus}}
\newcommand{\sigmoid}{\operatorname{sigmoid}}
\newcommand{\KL}{\operatorname{KL}}
\newcommand{\ELBO}{\operatorname{ELBO}}
\newcommand{\AIC}{\operatorname{AIC}}
\newcommand{\BIC}{\operatorname{BIC}}
\newcommand{\MAE}{\operatorname{MAE}}
\newcommand{\RMSE}{\operatorname{RMSE}}
\newcommand{\QLIKE}{\operatorname{QLIKE}}
\newcommand{\NLS}{\operatorname{NLS}}
\newcommand{\VaR}{\operatorname{VaR}}
\newcommand{\Data}{\mathcal{D}}
\newcommand{\Info}{\mathcal{F}}
\newcommand{\Param}{\symbf{\theta}}
\newcommand{\Xvec}{\symbf{X}}
\newcommand{\Pmat}{\symbf{P}}
\newcommand{\bmgamma}{\symbf{\gamma}}
\newcommand{\vecphi}{\symbf{\phi}}
\newcommand{\vech}{\operatorname{vech}}
\newcommand{\diff}{\mathrm{d}}

\title{\large\bfseries KHUNG LÝ THUYẾT VÀ QUY CÁCH TOÁN HỌC\\[0.35em]
\Large MÔ HÌNH MS-GARCH-X STUDENT-T CHO HỢP ĐỒNG TƯƠNG LAI VN30F1M:\\
SO SÁNH MÔ HÌNH VÀ SO SÁNH PHƯƠNG PHÁP SUY DIỄN}

\author{\textbf{Nguyễn Hoài Nam}\thanks{Khoa Toán -- Thống kê, Trường Đại học Tôn Đức Thắng. Email: \texttt{C2300152@student.tdtu.edu.vn}} \quad
\textbf{Nguyễn Thuận Hiếu}\footnotemark[1] \quad
\textbf{Đặng Chí Hiếu}\footnotemark[1] \\[0.3em]
\normalsize Giảng viên hướng dẫn: \textbf{TS. Trương Bửu Châu}\\[0.2em]
\normalsize Đề tài Nghiên cứu Khoa học Sinh viên Niên khóa 2026--2027}

\date{\normalsize Tháng 9 năm 2026}

\begin{document}
\maketitle
\thispagestyle{empty}

\begin{abstract}
\noindent Tài liệu thiết lập khung toán học cho đề tài mô hình hóa biến động hợp đồng
tương lai VN30F1M và giữ vai trò tài liệu quy chiếu về ký hiệu, đặc tả và tiêu chí đánh
giá cho toàn bộ nghiên cứu. Nghiên cứu được tổ chức thành hai phần độc lập về mục tiêu
nhưng dùng chung một khung đánh giá. Phần thứ nhất so sánh \emph{các mô hình} trong họ
GARCH bằng cùng một phương pháp ước lượng hợp lý cực đại, nhằm xác định giá trị gia tăng
của từng thành phần cấu trúc: đuôi dày, biến ngoại sinh và cơ chế chuyển trạng thái. Phần
thứ hai cố định mô hình chính MS-GARCH-X Student-t và so sánh \emph{các phương pháp suy
diễn}: hợp lý cực đại, biến phân mean-field và biến phân full-rank. Tài liệu đặc tả cơ
chế thu gọn moment kiểu Gray \cite{gray1996} cho lớp mô hình chuyển trạng thái, hàm hợp
lý qua bộ lọc tiến Hamilton \cite{hamilton1989}, khung Bayes biến phân với hiệu chỉnh
Jacobian trên không gian không ràng buộc, và hệ tiêu chí đánh giá ngoài mẫu dùng chung
cho cả hai phần.
\\[0.4em]
\textbf{Từ khóa:} MS-GARCH-X; Student-t chuẩn hóa; thu gọn moment Gray; bộ lọc Hamilton;
suy diễn biến phân; ELBO; nhận dạng nhãn trạng thái.
\end{abstract}

\vspace{0.2cm}

% ==============================================================================
\section*{Bảng ký hiệu}
\addcontentsline{toc}{section}{Bảng ký hiệu}
% ==============================================================================
\noindent Ký hiệu dưới đây thống nhất với bài báo chính của đề tài.

{\footnotesize
\begin{longtable}{@{}p{3.0cm}>{\raggedright\arraybackslash}p{12.2cm}@{}}
\toprule
\textbf{Ký hiệu} & \textbf{Ý nghĩa} \\
\midrule
\endfirsthead
\toprule
\textbf{Ký hiệu} & \textbf{Ý nghĩa} \\
\midrule
\endhead
\bottomrule
\endfoot
$r_t$ & Lợi suất logarit theo phần trăm tại phiên $t$. \\
$\varepsilon_t$ & Nhiễu của phương trình trung bình. \\
$z_t$ & Sai số chuẩn hóa, $\E[z_t]=0$ và $\Var(z_t)=1$. \\
$h_{t,s}$ & Phương sai có điều kiện tại $t$ khi $S_t=s$. \\
$\bar\mu_{t-1},\ \bar h_{t-1}$ & Trung bình và phương sai hỗn hợp thu gọn tại $t-1$. \\
$\bar\varepsilon_{t-1}$ & Phần dư thu gọn, $\bar\varepsilon_{t-1}=r_{t-1}-\bar\mu_{t-1}$. \\
$S_t$ & Trạng thái ẩn, $S_t\in\{1,2\}$. \\
$p_{ij}$ & Xác suất chuyển từ trạng thái $i$ sang $j$. \\
$\Pmat$ & Ma trận chuyển trạng thái Markov. \\
$\xi_{t\mid u}(s)$ & Xác suất $S_t=s$ có điều kiện trên thông tin đến $u$. \\
$\pi_s,\ d_s$ & Xác suất dừng và thời gian tồn tại kỳ vọng của trạng thái $s$. \\
$\mu_s,\omega_s,\alpha_s,\beta_s$ & Hệ số trung bình và phương sai của trạng thái $s$. \\
$\bmgamma_s$ & Vector hệ số biến ngoại sinh của trạng thái $s$. \\
$\Xvec_{t-1}$ & Vector $k$ biến ngoại sinh, biết trước thời điểm dự báo. \\
$\nu$ & Bậc tự do của phân phối Student-t chuẩn hóa. \\
$t^{*}_{\nu}$ & Phân phối Student-t chuẩn hóa về phương sai đơn vị. \\
$\Param$ & Vector toàn bộ tham số liên tục của mô hình. \\
$\vecphi$ & Ảnh của $\Param$ trên không gian không ràng buộc $\R^{d}$. \\
$\Info_{t-1}$ & Tập thông tin đến cuối phiên $t-1$. \\
$\Data_T$ & Toàn bộ dữ liệu dùng cho suy diễn đến $T$. \\
$\ell_T(\Param)$ & Log-hợp-lý quan sát trên mẫu huấn luyện. \\
$q(\vecphi)$ & Phân phối xấp xỉ hậu nghiệm trong suy diễn biến phân. \\
$RV_t$ & Biến động thực tế, đo trên thang phương sai. \\
$\VaR_{t,\alpha}$ & Giá trị chịu rủi ro tại mức xác suất đuôi $\alpha$. \\
\end{longtable}
}

% ==============================================================================
\section{Giới thiệu}\label{sec:intro}
% ==============================================================================

\subsection{Thiết kế nghiên cứu}\label{subsec:design}
Nghiên cứu được tổ chức thành hai phần, mỗi phần cố định một chiều so sánh và thay đổi
chiều còn lại. Cách tách này nhằm bảo đảm mỗi kết luận quy được về đúng một nguyên nhân.

\textbf{Phần thứ nhất --- so sánh mô hình.} Cố định phương pháp ước lượng là hợp lý cực
đại, thay đổi đặc tả mô hình theo một thứ bậc lồng nhau: từ GARCH chuẩn, thêm đuôi dày
Student-t, thêm biến ngoại sinh, rồi thêm cơ chế chuyển trạng thái. Vì các đặc tả chỉ
khác nhau ở một thành phần cấu trúc mỗi bước, chênh lệch kết quả giữa hai mô hình liền kề
đo được giá trị gia tăng của riêng thành phần vừa thêm.

\textbf{Phần thứ hai --- so sánh phương pháp suy diễn.} Cố định đặc tả là mô hình chính
MS-GARCH-X Student-t, thay đổi phương pháp suy diễn: hợp lý cực đại, biến phân mean-field
và biến phân full-rank. Vì cả ba làm việc trên cùng một hàm hợp lý và cùng một cơ chế thu
gọn moment, chênh lệch kết quả phản ánh khác biệt của bản thân phương pháp suy diễn chứ
không phải khác biệt của mô hình.

Hai phần dùng chung một khung đánh giá ngoài mẫu, trình bày ở \cref{sec:oos}, để kết quả
của chúng đặt cạnh nhau được.

\subsection{Phạm vi tài liệu}\label{subsec:scope}
Tài liệu đặc tả cơ sở toán học: cấu trúc mô hình, hàm hợp lý, khung suy diễn và tiêu chí
đánh giá. Tài liệu không trình bày chi tiết cài đặt tính toán --- lựa chọn thuật toán tối
ưu, siêu tham số điều khiển, cấu trúc chương trình --- vì những nội dung đó thuộc phần
thực nghiệm và không ảnh hưởng đến định nghĩa của các đại lượng được ước lượng. Khi một
lựa chọn tính toán \emph{có} hệ quả đối với cách diễn giải kết quả, hệ quả đó được nêu
thành nhận xét riêng.

% ==============================================================================
\section{Đặc tả mô hình}\label{sec:models}
% ==============================================================================

\subsection{Dữ liệu và tập thông tin}\label{subsec:data}
Xét thời gian rời rạc $t\in\{1,\dots,T\}$ và giá đóng cửa $P_t$. Lợi suất logarit theo
phần trăm:
\begin{equation}
  r_t = 100\times\ln\!\left(P_t/P_{t-1}\right).
  \label{eq:return}
\end{equation}
Vector $\Xvec_{t-1}\in\R^{k}$ gồm $k$ biến ngoại sinh trễ một bậc, chuẩn hóa z-score bằng
trung bình và độ lệch chuẩn \emph{của riêng tập huấn luyện}. Ràng buộc này là điều kiện
cần để loại trừ rò rỉ thông tin tương lai; dùng thống kê toàn mẫu sẽ làm kết quả ngoài
mẫu mất hiệu lực. Tập thông tin đến cuối phiên $t-1$ là
$\Info_{t-1}=\sigma\!\left(\{r_\tau\}_{\tau\le t-1},\{\Xvec_\tau\}_{\tau\le t-1}\right)$.
Đại lượng đại diện cho biến động thực tế là $RV_t$, đo trên thang phương sai, cùng đơn vị
với $h_{t,s}$.

\subsection{Trạng thái ẩn và xích Markov}\label{subsec:markov}
Chế độ thị trường là biến ẩn $S_t\in\{1,2\}$ tuân theo xích Markov thuần nhất bậc một:
\begin{equation}
  \Pmat=\begin{pmatrix} p_{11} & 1-p_{11}\\ 1-p_{22} & p_{22}\end{pmatrix},
  \qquad p_{ij}=\Prob\!\left(S_t=j\mid S_{t-1}=i\right).
  \label{eq:trans-mat}
\end{equation}
Khi $0<p_{11}<1$ và $0<p_{22}<1$, xích bất khả quy và không tuần hoàn, nên tồn tại duy
nhất phân phối dừng $\symbf\pi$ thỏa $\symbf\pi^{\top}\Pmat=\symbf\pi^{\top}$:
\begin{equation}
  \pi_1=\frac{1-p_{22}}{2-p_{11}-p_{22}},\qquad
  \pi_2=\frac{1-p_{11}}{2-p_{11}-p_{22}}.
  \label{eq:stationary-markov}
\end{equation}
Số phiên liên tiếp ở lại một trạng thái tuân theo phân phối hình học, nên thời gian tồn
tại kỳ vọng là
\begin{equation}
  d_s=\frac{1}{1-p_{ss}}.
  \label{eq:expected-duration}
\end{equation}
Đại lượng \eqref{eq:expected-duration} là cầu nối giữa tham số thống kê $p_{ss}$ và một
phát biểu kinh tế kiểm chứng được về độ bền của chế độ biến động.

\subsection{Phân phối sai số}\label{subsec:errors}
Mọi mô hình viết dưới dạng $r_t=\mu_{(\cdot)}+\varepsilon_t$ với
$\varepsilon_t=\sqrt{h_{t,(\cdot)}}\,z_t$ và $z_t$ độc lập cùng phân phối, chuẩn hóa sao
cho $\E[z_t]=0$, $\Var(z_t)=1$. Hai lựa chọn:
\begin{align}
  f_{\mathcal{N}}(z) &= \frac{1}{\sqrt{2\pi}}\exp\!\left(-\tfrac{1}{2}z^{2}\right),
  \label{eq:normal-density}\\[0.2em]
  f_{t^{*}}(z;\nu) &= \frac{\Gamma\!\left(\frac{\nu+1}{2}\right)}
  {\sqrt{\pi(\nu-2)}\,\Gamma\!\left(\frac{\nu}{2}\right)}
  \left(1+\frac{z^{2}}{\nu-2}\right)^{-\frac{\nu+1}{2}},\qquad \nu>2.
  \label{eq:std-t-density}
\end{align}
Mật độ \eqref{eq:std-t-density} là phân phối Student-t đã chuẩn hóa về phương sai đơn vị,
ký hiệu $t^{*}_{\nu}$. Thừa số $\nu-2$ thay cho $\nu$ dưới dấu căn chính là phép co giãn
bảo đảm $\Var(z_t)=1$; nhờ đó $h_{t,s}$ giữ đúng ý nghĩa phương sai có điều kiện và so
sánh được trực tiếp với trường hợp sai số chuẩn. Ràng buộc $\nu>2$ vì vậy không phải quy
ước kỹ thuật mà là điều kiện tồn tại của chính phép chuẩn hóa.

\subsection{Họ mô hình một trạng thái}\label{subsec:single-regime}
Ba mô hình đối chứng là các trường hợp riêng của một đặc tả chung:
\begin{equation}
  r_t=\mu+\varepsilon_t,\qquad
  h_t=\omega+\alpha\varepsilon_{t-1}^{2}+\beta h_{t-1}+\bmgamma^{\top}\Xvec_{t-1}.
  \label{eq:garchx}
\end{equation}

\begin{table}[H]
\centering\footnotesize
\setlength{\tabcolsep}{5pt}
\caption{Ba mô hình một trạng thái như trường hợp riêng của \eqref{eq:garchx}}
\label{tab:single-models}
\begin{tabularx}{\textwidth}{l c c c Y}
\toprule
\textbf{Mã} & \textbf{Phân phối $z_t$} & \textbf{$\bmgamma$} & \textbf{Số tham số} & \textbf{Thành phần cấu trúc được thêm} \\
\midrule
B1 & $\mathcal{N}(0,1)$ & $\equiv\symbf 0$ & $4$ & Mốc so sánh \\
B2 & $t^{*}_{\nu}$ & $\equiv\symbf 0$ & $5$ & Đuôi dày \\
B3 & $t^{*}_{\nu}$ & tự do & $5+k$ & Đuôi dày và biến ngoại sinh \\
\bottomrule
\end{tabularx}
\end{table}

\begin{assumption}[Điều kiện chính quy của mô hình một trạng thái]\label{as:single}
$\omega>0$, $\alpha\ge 0$, $\beta\ge 0$. Dừng hiệp phương sai yếu khi $\alpha+\beta<1$,
khi đó phương sai vô điều kiện là
\begin{equation}
  \bar h=\frac{\omega}{1-\alpha-\beta}.
  \label{eq:garch-unconditional-var}
\end{equation}
Với sai số $t^{*}_{\nu}$, moment bậc bốn tồn tại khi $\nu>4$ và
$3\alpha^{2}\frac{\nu-2}{\nu-4}+2\alpha\beta+\beta^{2}<1$.
\end{assumption}

\begin{remark}[Tính dương của phương sai khi có biến ngoại sinh]\label{rem:positivity}
Trong \eqref{eq:garchx}, tổ hợp $\bmgamma^{\top}\Xvec_{t-1}$ có thể nhận giá trị âm, nên
ràng buộc dấu của $(\omega,\alpha,\beta)$ \emph{không còn đủ} để bảo đảm $h_t>0$. Đây là
khác biệt bản chất giữa GARCH-X và GARCH cổ điển. Về mặt lý thuyết, tính dương phải được
xử lý bằng một cơ chế tường minh, hoặc bằng ràng buộc miền giá trị của
$\bmgamma^{\top}\Xvec_{t-1}$, hoặc bằng một hàm liên kết bảo toàn dấu. Trong nghiên cứu
này, tính dương được kiểm tra hậu nghiệm: tại nghiệm ước lượng của mọi mô hình, vế phải
của \eqref{eq:garchx} dương tại mọi thời điểm trong mẫu, nên đặc tả cộng tính giữ nguyên
hiệu lực và không cần đưa thêm hàm liên kết vào phương trình.
\end{remark}

\subsection{Phụ thuộc đường đi và thu gọn moment}\label{subsec:collapsing}
Khi thêm trạng thái ẩn, phương sai trở thành hàm của toàn bộ đường đi
$\symbf S_t=(S_1,\dots,S_t)$:
\begin{equation}
  h_t(\symbf S_t)=\omega_{S_t}+\alpha_{S_t}\varepsilon_{t-1}^{2}
  +\beta_{S_t}h_{t-1}(\symbf S_{t-1}).
  \label{eq:path-dependent-msgarch}
\end{equation}
Số nhánh tăng theo $2^{t}$, nên hàm hợp lý chính xác không tính được trong thời gian đa
thức. Gray \cite{gray1996} thay hai đại lượng trễ bằng moment hỗn hợp theo xác suất lọc
$\xi_{t-1\mid t-1}(i)$:
\begin{align}
  \bar\mu_{t-1} &= \sum_{i=1}^{2}\xi_{t-1\mid t-1}(i)\,\mu_i,
  \label{eq:collapsed-mean}\\
  \bar h_{t-1} &= \sum_{i=1}^{2}\xi_{t-1\mid t-1}(i)
  \left[h_{t-1,i}+\mu_i^{2}\right]-\bar\mu_{t-1}^{2},
  \label{eq:collapsed-variance}\\
  \bar\varepsilon_{t-1} &= r_{t-1}-\bar\mu_{t-1}.
  \label{eq:collapsed-residual}
\end{align}

\begin{proposition}[Dạng tương đương của \eqref{eq:collapsed-variance}]\label{prop:gray}
Phương trình \eqref{eq:collapsed-variance} tương đương với dạng phương sai của hỗn hợp
\begin{equation}
  \bar h_{t-1}=\sum_{i=1}^{2}\xi_{t-1\mid t-1}(i)
  \left[h_{t-1,i}+\left(\mu_i-\bar\mu_{t-1}\right)^{2}\right].
  \label{eq:collapsed-variance-alt}
\end{equation}
\end{proposition}
\begin{proof}
Đặt $\xi_i=\xi_{t-1\mid t-1}(i)$. Khai triển bình phương và dùng $\sum_i\xi_i=1$ cùng
\eqref{eq:collapsed-mean}:
$\sum_i\xi_i(\mu_i-\bar\mu_{t-1})^{2}
=\sum_i\xi_i\mu_i^{2}-2\bar\mu_{t-1}\sum_i\xi_i\mu_i+\bar\mu_{t-1}^{2}
=\sum_i\xi_i\mu_i^{2}-\bar\mu_{t-1}^{2}$.
Cộng $\sum_i\xi_i h_{t-1,i}$ vào hai vế thu được \eqref{eq:collapsed-variance}.
\end{proof}

Dạng \eqref{eq:collapsed-variance-alt} làm rõ ý nghĩa: phương sai hỗn hợp bằng trung bình
có trọng số của các phương sai trong từng trạng thái, cộng thêm độ phân tán của các trung
bình quanh trung bình hỗn hợp.

\begin{remark}[Vị trí của phép xấp xỉ]\label{rem:approx-location}
Sau khi thay $\bar\mu_{t-1}$ và $\bar h_{t-1}$ vào phương trình phương sai, quá trình trở
thành Markov và bộ lọc tiến ở \cref{subsec:filter} là \emph{chính xác} đối với mô hình đã
thu gọn: \eqref{eq:loglikelihood} là hàm hợp lý đúng của mô hình đó, không phải một xấp
xỉ của nó. Toàn bộ phép xấp xỉ nằm ở bước thu gọn
\eqref{eq:collapsed-mean}--\eqref{eq:collapsed-variance}. Hệ quả có tính phương pháp
luận và là điều kiện nền của Phần thứ hai: vì ba phương pháp suy diễn dùng chung một hàm
hợp lý, chênh lệch kết quả giữa chúng không quy cho khác biệt về cơ chế thu gọn được.
\end{remark}

\subsection{Mô hình chuyển trạng thái và mô hình chính}\label{subsec:main-model}
Mô hình chính của đề tài, ký hiệu M1:
\begin{align}
  r_t &= \mu_{S_t}+\varepsilon_t,
  \label{eq:main-mean}\\
  \varepsilon_t &= \sqrt{h_{t,S_t}}\,z_t,
  \label{eq:main-error}\\
  z_t &\sim t^{*}_{\nu},\qquad \nu>2,
  \label{eq:main-dist}\\
  h_{t,s} &= \omega_s+\alpha_s\bar\varepsilon_{t-1}^{2}
  +\beta_s\bar h_{t-1}+\bmgamma_s^{\top}\Xvec_{t-1},
  \qquad s\in\{1,2\}.
  \label{eq:main-variance}
\end{align}
Mô hình đối chứng B4 thu được bằng cách đặt $z_t\sim\mathcal{N}(0,1)$ và
$\bmgamma_s\equiv\symbf 0$; B4 tách riêng đóng góp của cơ chế chuyển trạng thái khỏi đóng
góp của đuôi dày và biến ngoại sinh. Vector tham số của M1:
\begin{equation}
  \Param=\left(\mu_1,\mu_2,\ \omega_1,\omega_2,\ \alpha_1,\alpha_2,\ \beta_1,\beta_2,\
  \bmgamma_1^{\top},\bmgamma_2^{\top},\ p_{11},p_{22},\ \nu\right)^{\top}
  \in\Theta\subset\R^{11+2k}.
  \label{eq:parameter-vector}
\end{equation}
Với $k=2$ biến ngoại sinh, $\dim(\Param)=15$. Bậc tự do $\nu$ dùng chung cho hai trạng
thái.

\begin{assumption}[Điều kiện chính quy của mô hình chuyển trạng thái]\label{as:ms}
\leavevmode
\begin{enumerate}
  \item[(C1)] \textbf{Tính ergodic:} $0<p_{11}<1$ và $0<p_{22}<1$, bảo đảm tồn tại duy
  nhất phân phối dừng \eqref{eq:stationary-markov}.
  \item[(C2)] \textbf{Tính dương của phương sai:} $\omega_s>0$, $\alpha_s\ge 0$,
  $\beta_s\ge 0$, và $h_{t,s}>0$ với mọi $t,s$.
  \item[(C3)] \textbf{Tồn tại phương sai có điều kiện:} $\nu>2$.
\end{enumerate}
\end{assumption}

\begin{remark}[Về điều kiện dừng bậc hai]\label{rem:stationarity}
Với lớp MS-GARCH của Haas, Mittnik và Paolella \cite{haas2004}, trong đó mỗi trạng thái
giữ một phương sai riêng không thu gọn, điều kiện dừng bậc hai có dạng đóng qua bán kính
phổ của một ma trận sinh từ $\Pmat$ và các hệ số $(\alpha_s,\beta_s)$. Biến thể thu gọn
kiểu Gray dùng ở đây \emph{không} thuộc lớp đó, nên điều kiện ấy không áp dụng trực tiếp
và tài liệu này không phát biểu nó như một tính chất của mô hình. Điểm cần ghi nhận là
ràng buộc $\alpha_s+\beta_s<1$ cho \emph{từng} trạng thái không phải điều kiện cần: một
trạng thái có thể bùng nổ cục bộ mà hệ thống vẫn dừng, miễn xác suất ở lại trạng thái đó
đủ nhỏ \cite{francq2001}. Vì vậy nghiên cứu áp đặt (C1)--(C3) như điều kiện vận hành và
đánh giá tính hợp lý của nghiệm qua $d_s$ theo \eqref{eq:expected-duration}.
\end{remark}

\subsection{Nhận dạng nhãn trạng thái}\label{subsec:label}
Hàm hợp lý bất biến dưới phép hoán vị nhãn: đổi đồng thời
$(\mu_1,\omega_1,\alpha_1,\beta_1,\bmgamma_1,p_{11})\leftrightarrow
(\mu_2,\omega_2,\alpha_2,\beta_2,\bmgamma_2,p_{22})$ cho giá trị hợp lý y hệt. Không gian
tham số vì thế chứa ít nhất hai điểm cực đại đối xứng, và mọi phát biểu về ``trạng thái
1'' chỉ có nghĩa sau khi cố định một quy ước.

\begin{definition}[Quy ước nhận dạng]\label{def:label}
Trạng thái được đánh số sao cho
\begin{equation}
  \E\!\left(h_t\mid S_t=1\right)<\E\!\left(h_t\mid S_t=2\right),
  \label{eq:label-identification}
\end{equation}
tức trạng thái $1$ là trạng thái biến động thấp.
\end{definition}

Quy ước áp dụng \emph{sau} khi ước lượng: cho nghiệm điểm trong khung hợp lý cực đại, và
cho \emph{từng mẫu} hậu nghiệm trong khung biến phân. Việc áp dụng theo từng mẫu là bắt
buộc: nếu chỉ sắp lại nhãn cho ước lượng tổng hợp, các moment hậu nghiệm đã bị trộn lẫn
giữa hai trạng thái trước đó và không khôi phục được.

\begin{remark}[Tỷ lệ hoán vị nhãn như một chẩn đoán]\label{rem:swap-rate}
Tỷ lệ mẫu hậu nghiệm phải hoán vị nhãn mang thông tin độc lập với các chẩn đoán hội tụ
thông thường. Tỷ lệ gần $0$ hoặc gần $1$ tương ứng với một quy ước nhất quán toàn cục và
vô hại. Tỷ lệ gần $1/2$ cho thấy phân phối hậu nghiệm không phân biệt được hai trạng
thái: danh tính trạng thái đổi qua lại ngay giữa các mẫu của cùng một lần chạy. Khi đó
mọi đại lượng theo trạng thái mất ý nghĩa, kể cả khi các chẩn đoán khác đạt chuẩn.
\end{remark}

% ==============================================================================
\section{Phần thứ nhất: so sánh mô hình bằng hợp lý cực đại}\label{sec:mle}
% ==============================================================================

Phần này cố định phương pháp ước lượng và so sánh năm đặc tả B1, B2, B3, B4 và M1.

\subsection{Bộ lọc tiến Hamilton và hàm hợp lý}\label{subsec:filter}
Với mô hình một trạng thái, log-hợp-lý tính trực tiếp từ mật độ có điều kiện. Với mô hình
chuyển trạng thái, biến ẩn $S_t$ được biên hóa bằng đệ quy tiến \cite{hamilton1989}, khởi
tạo $\xi_{1\mid 0}(s)=\pi_s$ theo \eqref{eq:stationary-markov}. Tại mỗi $t$:
\begin{align}
  \xi_{t\mid t-1}(j) &= \sum_{i=1}^{2}p_{ij}\,\xi_{t-1\mid t-1}(i),
  \label{eq:predicted-state-prob}\\
  \eta_{t,j} &= \frac{1}{\sqrt{h_{t,j}}}\,
  f\!\left(\frac{r_t-\mu_j}{\sqrt{h_{t,j}}}\right),
  \label{eq:state-density}\\
  f\!\left(r_t\mid\Info_{t-1};\Param\right)
  &= \sum_{j=1}^{2}\xi_{t\mid t-1}(j)\,\eta_{t,j},
  \label{eq:mixture-density}\\
  \xi_{t\mid t}(j) &= \frac{\xi_{t\mid t-1}(j)\,\eta_{t,j}}
  {f\!\left(r_t\mid\Info_{t-1};\Param\right)}.
  \label{eq:filtered-state-prob}
\end{align}
Bước \eqref{eq:filtered-state-prob} là công thức Bayes; thừa số $h_{t,j}^{-1/2}$ trong
\eqref{eq:state-density} là Jacobian của phép chuẩn hóa $z=(r_t-\mu_j)/\sqrt{h_{t,j}}$.
Log-hợp-lý quan sát:
\begin{equation}
  \ell_T(\Param)=\sum_{t=1}^{T}
  \ln\!\left[\sum_{j=1}^{2}\xi_{t\mid t-1}(j)\,\eta_{t,j}\right].
  \label{eq:loglikelihood}
\end{equation}
Khi cần mô tả lịch sử trạng thái, xác suất làm trơn tính bằng đệ quy lùi Kim
\cite{kim1994}:
\begin{equation}
  \xi_{t\mid T}(j)=\xi_{t\mid t}(j)\sum_{k=1}^{2}
  \frac{p_{jk}\,\xi_{t+1\mid T}(k)}{\xi_{t+1\mid t}(k)}.
  \label{eq:kim-smoother}
\end{equation}
Đại lượng \eqref{eq:kim-smoother} sử dụng thông tin sau thời điểm $t$ nên chỉ dùng để mô
tả lịch sử, không dùng làm tín hiệu thời gian thực.

\subsection{Tham số hóa không ràng buộc}\label{subsec:reparam}
Bài toán $\max_{\Param\in\Theta}\ell_T(\Param)$ có ràng buộc. Ta chuyển sang biến không
ràng buộc $\vecphi\in\R^{d}$ qua song ánh trơn $\Param=g(\vecphi)$, tác động theo từng
tọa độ:
\begin{equation}
  \mu_s=\phi_{\mu,s},\quad
  \bmgamma_s=\vecphi_{\gamma,s},\quad
  \omega_s=e^{\phi_{\omega,s}},\quad
  \alpha_s=\softplus(\phi_{\alpha,s}),\quad
  \beta_s=\softplus(\phi_{\beta,s}),
  \label{eq:transform-positive}
\end{equation}
\begin{equation}
  p_{ss}=\sigmoid(\phi_{p,s}),\qquad
  \nu=2+\softplus(\phi_{\nu}).
  \label{eq:transform-bounded}
\end{equation}
Phép biến đổi này bảo đảm (C1)--(C3) theo cấu trúc chứ không bằng kiểm tra hậu nghiệm:
ảnh của $\softplus$ là $(0,\infty)$, ảnh của $\sigmoid$ là $(0,1)$, và
$\nu=2+\softplus(\phi_\nu)>2$. Lựa chọn $\omega_s=e^{\phi}$ tương thích với tiên nghiệm
log-chuẩn ở \cref{subsec:priors}, bảo đảm hai khung suy diễn của Phần thứ hai làm việc
trên cùng một tham số hóa.

\subsection{Tính đa cực trị và nguyên tắc chọn nghiệm}\label{subsec:multimodality}
Hàm hợp lý của mô hình chuyển trạng thái đa cực trị. Điều này không phải khiếm khuyết của
thuật toán mà là tính chất của chính hàm mục tiêu, bắt nguồn từ hai nguồn: đối xứng hoán
vị nhãn đã nêu ở \cref{subsec:label}, và sự tồn tại của các cực đại giả.

\begin{remark}[Cực đại giả và điều kiện chính quy vận hành]\label{rem:spurious}
Hàm hợp lý MS-GARCH có các cực đại với giá trị tùy ý lớn, ứng với một trạng thái tồn tại
đúng một phiên và hấp thụ trọn một quan sát cực đoan: phương sai của trạng thái đó co về
gần không và mật độ tại điểm quan sát phân kỳ. Những điểm này tối ưu về mặt số học nhưng
vô nghĩa về mặt kinh tế, và một quy trình chọn nghiệm thuần túy theo giá trị hợp lý sẽ
chọn đúng chúng. Nghiên cứu vì vậy chọn nghiệm có hợp lý cao nhất \emph{trong số các
nghiệm thỏa điều kiện chính quy vận hành}
\begin{equation}
  d_s\ge 2\ \text{phiên}\quad\text{và}\quad p_{ss}\in[0{,}01;\ 0{,}999],
  \qquad s=1,2,
  \label{eq:regularity}
\end{equation}
tức loại bỏ các nghiệm mô tả một chế độ không tồn tại đủ lâu để mang ý nghĩa chế độ, hoặc
một chế độ gần như hấp thụ. Quy tắc này phải được công bố cùng kết quả, vì nó có thể loại
bỏ những nghiệm có log-hợp-lý cao hơn nghiệm được báo cáo.
\end{remark}

Vì hàm hợp lý đa cực trị, ước lượng được thực hiện từ nhiều điểm khởi tạo độc lập, trong
đó một phần sinh từ nghiệm của mô hình lồng nhau đơn giản hơn theo thứ bậc B1 $\to$ B4 và
B3 $\to$ M1. Độ tán của giá trị hợp lý giữa các điểm khởi tạo là một chẩn đoán có giá
trị độc lập: độ tán nhỏ củng cố tin cậy vào nghiệm, độ tán lớn cho thấy bề mặt hợp lý
nhiều đỉnh và kết luận cần thận trọng tương ứng.

\subsection{Sai số chuẩn và điều kiện hiệu lực}\label{subsec:se}
Gọi $\hat{\vecphi}$ là nghiệm và
$\symbf H(\hat{\vecphi})=\nabla^{2}_{\vecphi}\ell_T(g(\vecphi))\big|_{\hat{\vecphi}}$ là
Hessian quan sát. Khi $\symbf H(\hat{\vecphi})$ xác định âm, hiệp phương sai tiệm cận
trên không gian không ràng buộc là $[-\symbf H(\hat{\vecphi})]^{-1}$, và sai số chuẩn của
tham số gốc suy ra bằng phương pháp delta:
\begin{equation}
  \widehat{\Cov}\!\left(\hat{\Param}\right)
  =\symbf J_g(\hat{\vecphi})\left[-\symbf H(\hat{\vecphi})\right]^{-1}
  \symbf J_g(\hat{\vecphi})^{\top},
  \qquad
  \symbf J_g(\vecphi)=\frac{\partial g(\vecphi)}{\partial\vecphi^{\top}}.
  \label{eq:delta-method}
\end{equation}

\begin{remark}[Điều kiện hiệu lực của \eqref{eq:delta-method}]\label{rem:se-validity}
Tính chuẩn tiệm cận của ước lượng hợp lý cực đại giả định nghiệm nằm \emph{bên trong}
không gian tham số. Khi nghiệm chạm hoặc gần chạm biên --- cụ thể khi $p_{ss}$ tiến sát
$1$, tức chế độ gần như hấp thụ --- giả định bị vi phạm và \eqref{eq:delta-method} không
còn là sai số chuẩn hợp lệ; đại lượng thu được chỉ đo độ cong cục bộ của hàm hợp lý. Khi
điều này xảy ra, kết quả phải được báo cáo kèm cảnh báo, và các phát biểu về ý nghĩa
thống kê của tham số tương ứng phải được rút lại.
\end{remark}

\subsection{Tiêu chí so sánh mô hình}\label{subsec:ic}
Với $d=\dim(\Param)$ và $T$ là cỡ mẫu huấn luyện:
\begin{equation}
  \AIC=2d-2\ell_T,\qquad \BIC=d\ln T-2\ell_T.
  \label{eq:information-criteria}
\end{equation}
Cả hai phạt độ phức tạp, nhưng $\BIC$ phạt nặng hơn khi $T>e^{2}\approx 7{,}4$. Với $T$
hàng nghìn, khoảng cách giữa hai tiêu chí đủ lớn để chúng xếp hạng khác nhau; khi đó
$\BIC$ nghiêng về mô hình tiết kiệm tham số hơn. Trong thiết kế lồng nhau của Phần thứ
nhất, chênh lệch $\AIC$ và $\BIC$ giữa hai mô hình liền kề là thước đo trực tiếp giá trị
gia tăng của thành phần cấu trúc vừa thêm, sau khi đã trừ chi phí tham số.

\begin{remark}[Vì sao không dùng kiểm định tỷ số hợp lý]\label{rem:lr-test}
So sánh mô hình một trạng thái với mô hình hai trạng thái không thực hiện được bằng kiểm
định tỷ số hợp lý. Dưới giả thuyết không ``chỉ có một trạng thái'', các tham số
$p_{11},p_{22}$ và bộ tham số của trạng thái thứ hai \emph{không được nhận dạng}: hàm hợp
lý không phụ thuộc vào chúng. Điều kiện chính quy của định lý Wilks bị vi phạm và thống
kê tỷ số hợp lý không có phân phối tiệm cận $\chi^{2}$ \cite{davies1987,hansen1992}. Đây
là lý do thứ bậc mô hình được xếp bằng \eqref{eq:information-criteria} và bằng năng lực
dự báo ngoài mẫu, chứ không bằng kiểm định giả thuyết.
\end{remark}

% ==============================================================================
\section{Phần thứ hai: so sánh phương pháp suy diễn}\label{sec:vb}
% ==============================================================================

Phần này cố định đặc tả là mô hình chính M1 và so sánh ba phương pháp suy diễn. Hợp lý
cực đại đã trình bày ở \cref{sec:mle}; mục này xây dựng khung Bayes biến phân.

\subsection{Tiên nghiệm}\label{subsec:priors}
Giả định độc lập tiên nghiệm giữa các khối tham số, $p(\Param)=\prod_i p(\Param_i)$:
\begin{equation}
\begin{aligned}
  \mu_s &\sim \mathcal{N}(0,\ \sigma_\mu^{2}), &
  \ln\omega_s &\sim \mathcal{N}(m_\omega,\ \sigma_\omega^{2}), \\
  \alpha_s &\sim \operatorname{Beta}(a,\ a), &
  \beta_s &\sim \operatorname{Beta}(a,\ a), \\
  \bmgamma_s &\sim \mathcal{N}(\symbf 0,\ \sigma_\gamma^{2}\symbf I_k), &
  p_{ss} &\sim \operatorname{Beta}(a_p,\ b_p), \\
  \nu-2 &\sim \operatorname{Gamma}(a_\nu,\ b_\nu). & &
\end{aligned}
\label{eq:priors}
\end{equation}
Ba lựa chọn mang nội dung thực chất. Tiên nghiệm log-chuẩn trên $\omega_s$ tương thích
với tham số hóa \eqref{eq:transform-positive}, bảo đảm hai khung suy diễn làm việc trên
cùng một miền. Tiên nghiệm $\operatorname{Beta}$ trên $p_{ss}$ với $a_p>b_p$ đặt phần lớn
khối lượng ở vùng $p_{ss}$ lớn, phản ánh nhận định rằng các chế độ biến động có tính dai
dẳng. Tiên nghiệm $\operatorname{Gamma}$ đặt trên $\nu-2$ thay vì trên $\nu$ bảo đảm
$\nu>2$ theo cấu trúc, nhất quán với (C3).

\begin{remark}[Hệ quả của tiên nghiệm Beta đối với việc so sánh]\label{rem:prior-support}
Tiên nghiệm $\operatorname{Beta}$ chặn cứng $\alpha_s,\beta_s$ trong $(0,1)$, trong khi
tham số hóa softplus \eqref{eq:transform-positive} của hợp lý cực đại không có chặn trên.
Hai khung suy diễn vì vậy làm việc trên hai miền tham số khác nhau, và chênh lệch ước
lượng $\beta_s$ giữa chúng không quy hoàn toàn cho khác biệt về nguyên lý suy diễn được.
Đây là một giới hạn cấu trúc của thiết kế so sánh trong Phần thứ hai, phải được nêu khi
diễn giải kết quả.
\end{remark}

\subsection{Hậu nghiệm trên không gian không ràng buộc}\label{subsec:posterior}
Theo định lý Bayes, $p(\Param\mid\Data_T)\propto p(\Param)\,\mathcal{L}(\Param;\Data_T)$
với $\mathcal{L}(\Param;\Data_T)=\exp\ell_T(\Param)$ theo \eqref{eq:loglikelihood}. Suy
diễn thực hiện trên $\vecphi$ với cùng phép biến đổi
\eqref{eq:transform-positive}--\eqref{eq:transform-bounded}. Theo công thức đổi biến:
\begin{equation}
  p(\vecphi\mid\Data_T)\propto
  p\!\left(g(\vecphi)\right)\cdot\mathcal{L}\!\left(g(\vecphi);\Data_T\right)
  \cdot\left|\det\symbf J_g(\vecphi)\right|.
  \label{eq:posterior-unconstrained}
\end{equation}
Vì $g$ tác động theo từng tọa độ, $\symbf J_g$ là ma trận đường chéo và
\begin{equation}
  \ln\left|\det\symbf J_g(\vecphi)\right|
  =\sum_{i=1}^{d}\ln\left|\frac{\partial g_i(\phi_i)}{\partial\phi_i}\right|.
  \label{eq:log-jacobian}
\end{equation}
Số hạng \eqref{eq:log-jacobian} không phải chi tiết kỹ thuật có thể bỏ qua: thiếu nó,
mật độ hậu nghiệm trên $\vecphi$ không còn là ảnh đúng của hậu nghiệm trên $\Theta$, và
mọi moment tính từ đó đều lệch một cách hệ thống.

\subsection{Cận dưới bằng chứng}\label{subsec:elbo}
Tích phân chuẩn hóa $p(\Data_T)$ không có dạng đóng. Suy diễn biến phân chọn một họ phân
phối $\mathcal{Q}$ và tìm phần tử gần hậu nghiệm nhất theo phân kỳ Kullback--Leibler
\cite{blei2017}. Với mọi $q\in\mathcal{Q}$ ta có phân rã đồng nhất
\begin{equation}
  \ln p(\Data_T)=\ELBO(q)+\KL\!\left(q(\vecphi)\,\|\,p(\vecphi\mid\Data_T)\right).
  \label{eq:elbo-decomposition}
\end{equation}
Vì $\ln p(\Data_T)$ không phụ thuộc $q$ và $\KL\ge 0$, cực tiểu hóa $\KL$ tương đương cực
đại hóa $\ELBO$, và $\ELBO(q)\le\ln p(\Data_T)$ với dấu bằng khi và chỉ khi $q$ trùng hậu
nghiệm. Dạng tường minh:
\begin{equation}
  \ELBO(q)=\E_{q}\!\left[
  \ell_T\!\left(g(\vecphi)\right)+\ln p\!\left(g(\vecphi)\right)
  +\ln\left|\det\symbf J_g(\vecphi)\right|\right]+\mathcal{H}(q),
  \label{eq:elbo}
\end{equation}
với $\mathcal{H}(q)=-\E_q[\ln q(\vecphi)]$ là entropy vi phân của phân phối biến phân.

\begin{remark}[Phạm vi so sánh được của $\ELBO$]\label{rem:elbo-scope}
$\ELBO$ là cận dưới của $\ln p(\Data_T)$ \emph{dưới một mô hình xác suất cụ thể}. Khi đổi
tiên nghiệm, đại lượng được cận dưới cũng đổi theo, nên $\ELBO$ của hai biến thể tiên
nghiệm khác nhau không so sánh trực tiếp được và không dùng để chọn tiên nghiệm. Trong
cùng một đặc tả, $\ELBO$ so sánh được giữa mean-field và full-rank, vì khi đó chỉ họ
$\mathcal{Q}$ thay đổi --- và đây chính là cơ sở để đối chiếu hai phương pháp biến phân
trong Phần thứ hai.
\end{remark}

\subsection{Hai họ phân phối biến phân}\label{subsec:families}

\begin{definition}[Mean-field]\label{def:mf}
Họ Gauss thừa số hóa theo tọa độ:
\begin{equation}
  q_{\mathrm{MF}}(\vecphi)=\prod_{i=1}^{d}\mathcal{N}\!\left(\phi_i\mid m_i,s_i^{2}\right)
  =\mathcal{N}\!\left(\vecphi\mid\symbf m,\ \diag(\symbf s^{2})\right),
  \label{eq:mean-field}
\end{equation}
với tham số biến phân $(\symbf m,\symbf s)$ gồm $2d$ thành phần; entropy có dạng đóng
$\mathcal{H}(q_{\mathrm{MF}})=\frac{d}{2}(1+\ln 2\pi)+\sum_i\ln s_i$; gradient của
\eqref{eq:elbo} ước lượng không chệch qua tái tham số hóa $\phi_i=m_i+s_i\varepsilon_i$
với $\varepsilon_i\sim\mathcal{N}(0,1)$ độc lập \cite{kingma2014}.
\end{definition}

\begin{definition}[Full-rank]\label{def:fr}
Họ Gauss đa biến với hiệp phương sai đầy đủ qua thừa số Cholesky:
\begin{equation}
  q_{\mathrm{FR}}(\vecphi)=\mathcal{N}\!\left(\vecphi\mid\symbf m,\ \symbf\Sigma\right),
  \qquad \symbf\Sigma=\symbf L\symbf L^{\top},
  \label{eq:full-rank}
\end{equation}
với $\symbf L$ tam giác dưới có đường chéo dương; tham số biến phân
$(\symbf m,\vech(\symbf L))$ gồm $d+d(d+1)/2$ thành phần; entropy
$\mathcal{H}(q_{\mathrm{FR}})=\frac{d}{2}(1+\ln 2\pi)+\sum_i\ln L_{ii}$; tái tham số hóa
$\vecphi=\symbf m+\symbf L\symbf\varepsilon$ với
$\symbf\varepsilon\sim\mathcal{N}(\symbf 0,\symbf I_d)$.
\end{definition}

Phân rã $\symbf\Sigma=\symbf L\symbf L^{\top}$ với $L_{ii}>0$ bảo đảm $\symbf\Sigma$ xác
định dương theo cấu trúc và là duy nhất; đây là lý do tối ưu hóa thực hiện trên
$\symbf L$ chứ không trực tiếp trên $\symbf\Sigma$.

\begin{remark}[Đánh đổi giữa hai họ]\label{rem:mf-vs-fr}
Với $d=15$, mean-field tối ưu $30$ tham số còn full-rank tối ưu $135$. Mean-field áp đặt
$\Cov_q(\phi_i,\phi_j)=0$, triệt tiêu tương quan hậu nghiệm --- đáng kể nhất là tương quan
âm giữa $\alpha_s$ và $\beta_s$, vốn là đặc trưng của họ GARCH --- và do dạng bất đối
xứng của $\KL(q\|p)$, xu hướng hệ thống là đánh giá thấp phương sai hậu nghiệm.
Full-rank không có hạn chế đó và về lý thuyết cho cận dưới khít hơn. Tuy nhiên, vì cả hai
chỉ cho \emph{cận dưới}, một họ giàu hơn không bảo đảm nghiệm tốt hơn trên thực tế: khi
bài toán tối ưu $135$ tham số không được giải tốt, full-rank có thể cho $\ELBO$ thấp hơn
mean-field. Đây là khả năng phải kiểm tra bằng thực nghiệm, không loại trừ được bằng lý
thuyết, và là một trong những câu hỏi mà Phần thứ hai trả lời.
\end{remark}

\subsection{Đối chiếu ba phương pháp}\label{subsec:comparison}

\begin{table}[H]
\centering\footnotesize
\setlength{\tabcolsep}{4pt}
\caption{Đối chiếu ba phương pháp suy diễn trên mô hình MS-GARCH-X Student-t}
\label{tab:comparison}
\begin{tabularx}{\textwidth}{>{\raggedright\arraybackslash\bfseries}p{2.5cm} Y Y Y}
\toprule
\textbf{Tiêu chí} & \textbf{Hợp lý cực đại} & \textbf{Biến phân mean-field} & \textbf{Biến phân full-rank} \\
\midrule
Đối tượng tối ưu & $\ell_T(\Param)$, \eqref{eq:loglikelihood} & $\ELBO(q)$, \eqref{eq:elbo} & $\ELBO(q)$, \eqref{eq:elbo} \\
\addlinespace[0.2em]
Kết quả suy diễn & Một điểm $\hat{\Param}$ & Phân phối $q_{\mathrm{MF}}$ & Phân phối $q_{\mathrm{FR}}$ \\
\addlinespace[0.2em]
Vai trò tiên nghiệm & Không có & \eqref{eq:priors} kèm \eqref{eq:log-jacobian} & \eqref{eq:priors} kèm \eqref{eq:log-jacobian} \\
\addlinespace[0.2em]
Hiệp phương sai & Tiệm cận qua $[-\symbf H]^{-1}$, hiệu lực có điều kiện & Đường chéo theo giả định & Đầy đủ qua $\symbf L\symbf L^{\top}$ \\
\addlinespace[0.2em]
Số tham số tối ưu & $d$ & $2d$ & $d+d(d+1)/2$ \\
\addlinespace[0.2em]
Dự báo mật độ & Cắm vào $\hat{\Param}$ & \eqref{eq:posterior-predictive} & \eqref{eq:posterior-predictive} \\
\addlinespace[0.2em]
Hạn chế chính & Suy diễn tiệm cận mất hiệu lực khi nghiệm sát biên & Bỏ tương quan; đánh giá thấp bất định & Bài toán tối ưu khó, không bảo đảm cận tốt hơn \\
\bottomrule
\end{tabularx}
\end{table}

\begin{remark}[Điều kiện để so sánh có hiệu lực]\label{rem:comparability}
Ba phương pháp chỉ so sánh được khi dùng chung hàm hợp lý, chung cơ chế thu gọn và chung
quy ước nhận dạng nhãn. Hai điều kiện đầu được bảo đảm vì cả ba cùng dựa trên
\eqref{eq:loglikelihood} với cùng bước thu gọn
\eqref{eq:collapsed-mean}--\eqref{eq:collapsed-variance} (\cref{rem:approx-location});
điều kiện thứ ba được bảo đảm bằng \cref{def:label} áp dụng nhất quán. Điểm bất đối xứng
còn lại --- hợp lý cực đại không có tiên nghiệm trong khi hai phương pháp biến phân có ---
là khác biệt bản chất giữa hai khung suy diễn, không loại bỏ được, và phải được tính đến
khi diễn giải mọi chênh lệch ước lượng.
\end{remark}

% ==============================================================================
\section{Khung đánh giá ngoài mẫu}\label{sec:oos}
% ==============================================================================

Khung này dùng chung cho cả hai phần. Mẫu được chia theo thời gian: $T$ quan sát đầu để
huấn luyện, $N$ quan sát cuối để kiểm định. Không tham số nào được ước lượng lại bằng dữ
liệu kiểm định.

\subsection{Dự báo phương sai}\label{subsec:vol-forecast}
Với $\widehat h_t$ là phương sai dự báo một bước và $RV_t$ là đại lượng đại diện, cả hai
trên thang phương sai:
\begin{align}
  \MAE &= \frac{1}{N}\sum_{t}\left|RV_t-\widehat h_t\right|,
  \label{eq:mae-vol}\\
  \RMSE &= \sqrt{\frac{1}{N}\sum_{t}\left(RV_t-\widehat h_t\right)^{2}},
  \label{eq:rmse-vol}\\
  \QLIKE &= \frac{1}{N}\sum_{t}\left[\ln\widehat h_t+\frac{RV_t}{\widehat h_t}\right].
  \label{eq:qlike}
\end{align}

\begin{remark}[Vì sao $\QLIKE$ là tiêu chí chính]\label{rem:qlike}
$RV_t$ chỉ là một đại lượng \emph{đại diện} có nhiễu cho phương sai thật, vốn không quan
sát được. Theo Patton \cite{patton2011}, $\QLIKE$ thuộc lớp hàm tổn thất \emph{vững} đối
với loại nhiễu này: thứ hạng mô hình theo $\QLIKE$ kỳ vọng không đổi khi thay phương sai
thật bằng một đại lượng đại diện không chệch. $\MAE$ và $\RMSE$ không có tính chất này ---
chúng có thể xếp hạng sai khi đại lượng đại diện nhiễu --- nên được báo cáo bổ sung chứ
không dùng làm căn cứ kết luận.
\end{remark}

\subsection{Dự báo mật độ}\label{subsec:density-forecast}
Dự báo mật độ đánh giá toàn bộ phân phối dự báo $\widehat f_t$, không chỉ moment bậc hai:
\begin{equation}
  \NLS=-\frac{1}{N}\sum_{t}\ln\widehat f_t(r_t).
  \label{eq:negative-log-score}
\end{equation}
$\NLS$ là một quy tắc chấm điểm chặt: kỳ vọng của nó đạt cực tiểu khi và chỉ khi phân
phối báo cáo trùng phân phối sinh dữ liệu. Đây là căn cứ để dùng $\NLS$ so sánh các mô
hình khác nhau về hình dạng đuôi, điều mà $\QLIKE$ không làm được vì chỉ nhìn moment bậc
hai.

Cơ chế sinh $\widehat f_t$ khác nhau giữa hai khung suy diễn, và khác biệt này là một
trong những đối tượng so sánh của Phần thứ hai. Hợp lý cực đại dùng ước lượng cắm vào,
$\widehat f_t(\cdot)=f(\cdot\mid\Info_{t-1};\hat{\Param})$, bỏ qua bất định tham số. Các
phương pháp biến phân tích phân trên hậu nghiệm:
\begin{equation}
  \widehat f_t(\cdot)
  =\int_{\Theta}f\!\left(\cdot\mid\Info_{t-1};\Param\right)p(\Param\mid\Data_T)\,\diff\Param
  \approx\frac{1}{S}\sum_{s=1}^{S}
  f\!\left(\cdot\mid\Info_{t-1};g(\vecphi^{(s)})\right),
  \qquad \vecphi^{(s)}\sim q.
  \label{eq:posterior-predictive}
\end{equation}
Phân phối \eqref{eq:posterior-predictive} là hỗn hợp của $S$ mật độ nên có đuôi dày hơn
từng thành phần. Đây là cơ chế qua đó bất định tham số đi vào dự báo, và là lý do các
phương pháp Bayes có lợi thế tiềm năng ở $\NLS$ ngay cả khi không có lợi thế ở $\QLIKE$.

\subsection{Kiểm định giá trị chịu rủi ro}\label{subsec:var-tests}
Gọi $I_t=\symbf 1\{r_t<\widehat{\VaR}_{t,\alpha}\}$ là chỉ báo vi phạm. Một mô hình rủi
ro đúng đòi hỏi hai tính chất: tỷ lệ vi phạm bằng $\alpha$, và các vi phạm độc lập theo
thời gian.
\begin{itemize}
  \item \textbf{Kupiec} \cite{kupiec1995} kiểm định độ phủ vô điều kiện
  $H_0\!:\E[I_t]=\alpha$, với thống kê $LR_{uc}\overset{d}{\to}\chi^{2}(1)$.
  \item \textbf{Christoffersen} \cite{christoffersen1998} bổ sung kiểm định tính độc lập
  $LR_{ind}\overset{d}{\to}\chi^{2}(1)$ trên chuỗi $\{I_t\}$ qua mô hình Markov hai trạng
  thái. Thống kê độ phủ có điều kiện là $LR_{cc}=LR_{uc}+LR_{ind}$, với
  $LR_{cc}\overset{d}{\to}\chi^{2}(2)$.
\end{itemize}
Việc tách hai thành phần có ý nghĩa thực chất: một mô hình có đúng số vi phạm nhưng dồn
toàn bộ vào một giai đoạn căng thẳng sẽ vượt qua Kupiec nhưng bị $LR_{ind}$ bác bỏ, và đó
đúng là kiểu sai hỏng nguy hiểm nhất trong quản trị rủi ro. Kiểm định VaR vì vậy là nơi
giả định đuôi dày được kiểm chứng trực tiếp, chứ không chỉ được biện luận.

So sánh cặp mô hình theo hàm tổn thất được thực hiện bằng kiểm định Diebold--Mariano
\cite{diebold1995} với hiệu chỉnh tự tương quan; khi số cặp so sánh lớn, kết luận cần
tính đến vấn đề kiểm định bội.

% ==============================================================================
\section{Tổng hợp thiết kế thực nghiệm}\label{sec:design-summary}
% ==============================================================================

\begin{table}[H]
\centering\footnotesize
\setlength{\tabcolsep}{4pt}
\caption{Hai phần của nghiên cứu: đối tượng so sánh, đại lượng cố định và tiêu chí}
\label{tab:design}
\begin{tabularx}{\textwidth}{>{\raggedright\arraybackslash\bfseries}p{2.2cm} Y Y}
\toprule
& \textbf{Phần thứ nhất} & \textbf{Phần thứ hai} \\
\midrule
Câu hỏi & Thành phần cấu trúc nào cải thiện mô hình hóa biến động? & Phương pháp suy diễn nào phù hợp cho mô hình chính? \\
\addlinespace[0.2em]
Cố định & Phương pháp: hợp lý cực đại & Mô hình: MS-GARCH-X Student-t (M1) \\
\addlinespace[0.2em]
Thay đổi & Đặc tả: B1, B2, B3, B4, M1 & Phương pháp: hợp lý cực đại, biến phân mean-field, biến phân full-rank \\
\addlinespace[0.2em]
Tiêu chí trong mẫu & $\AIC$, $\BIC$ theo \eqref{eq:information-criteria} & $\ELBO$ theo \eqref{eq:elbo}; so sánh giữa hai họ biến phân \\
\addlinespace[0.2em]
Tiêu chí ngoài mẫu & $\QLIKE$, $\NLS$, kiểm định VaR & $\QLIKE$, $\NLS$, kiểm định VaR; độ rộng khoảng dự báo \\
\addlinespace[0.2em]
Điều kiện hiệu lực & Cùng dữ liệu, cùng phân chia mẫu, cùng quy ước nhãn & Cùng hàm hợp lý, cùng cơ chế thu gọn, cùng quy ước nhãn \\
\bottomrule
\end{tabularx}
\end{table}

Hai phần bổ sung cho nhau về mặt kết luận. Phần thứ nhất trả lời câu hỏi \emph{nên dùng
mô hình nào}; Phần thứ hai trả lời câu hỏi \emph{với mô hình đó, nên ước lượng bằng cách
nào}, và đồng thời cung cấp một đánh giá về mức độ mà kết luận của Phần thứ nhất phụ
thuộc vào phương pháp ước lượng được chọn.

% ==============================================================================
\begin{thebibliography}{99}
\bibitem{blei2017} Blei, D. M., Kucukelbir, A., \& McAuliffe, J. D. (2017). Variational inference: A review for statisticians. \emph{Journal of the American Statistical Association}, 112(518), 859--877.
\bibitem{bollerslev1986} Bollerslev, T. (1986). Generalized autoregressive conditional heteroskedasticity. \emph{Journal of Econometrics}, 31(3), 307--327.
\bibitem{christoffersen1998} Christoffersen, P. F. (1998). Evaluating interval forecasts. \emph{International Economic Review}, 39(4), 841--862.
\bibitem{davies1987} Davies, R. B. (1987). Hypothesis testing when a nuisance parameter is present only under the alternative. \emph{Biometrika}, 74(1), 33--43.
\bibitem{diebold1995} Diebold, F. X., \& Mariano, R. S. (1995). Comparing predictive accuracy. \emph{Journal of Business \& Economic Statistics}, 13(3), 253--263.
\bibitem{francq2001} Francq, C., \& Zakoïan, J. M. (2001). Stationarity of multivariate Markov-switching ARMA models. \emph{Journal of Econometrics}, 102(2), 339--364.
\bibitem{gray1996} Gray, S. F. (1996). Modeling the conditional distribution of interest rates as a regime-switching process. \emph{Journal of Financial Economics}, 42(1), 27--62.
\bibitem{haas2004} Haas, M., Mittnik, S., \& Paolella, M. S. (2004). A new approach to Markov-switching GARCH models. \emph{Journal of Financial Econometrics}, 2(4), 493--530.
\bibitem{hamilton1989} Hamilton, J. D. (1989). A new approach to the economic analysis of nonstationary time series and the business cycle. \emph{Econometrica}, 57(2), 357--384.
\bibitem{hansen1992} Hansen, B. E. (1992). The likelihood ratio test under nonstandard conditions: Testing the Markov switching model of GNP. \emph{Journal of Applied Econometrics}, 7(S1), S61--S82.
\bibitem{kim1994} Kim, C. J. (1994). Dynamic linear models with Markov-switching. \emph{Journal of Econometrics}, 60(1--2), 1--22.
\bibitem{kingma2014} Kingma, D. P., \& Welling, M. (2014). Auto-encoding variational Bayes. \emph{International Conference on Learning Representations}.
\bibitem{klaassen2002} Klaassen, F. (2002). Improving GARCH volatility forecasts with regime-switching GARCH. \emph{Empirical Economics}, 27(2), 363--394.
\bibitem{kupiec1995} Kupiec, P. H. (1995). Techniques for verifying the accuracy of risk measurement models. \emph{The Journal of Derivatives}, 3(2), 73--84.
\bibitem{nelson1990} Nelson, D. B. (1990). Stationarity and persistence in GARCH(1,1) models. \emph{Econometric Theory}, 6(3), 318--334.
\bibitem{patton2011} Patton, A. J. (2011). Volatility forecast comparison using imperfect volatility proxies. \emph{Journal of Econometrics}, 160(1), 246--256.
\end{thebibliography}

\end{document}
